% !TeX encoding = UTF-8
% !TeX program = pdflatex
% IEEEtran Journal Template -- Migrated from Elsevier cja-template.tex
\documentclass[lettersize,journal]{IEEEtran}

%% ---- Packages (only those compatible with IEEEtran) ----
\usepackage{amsmath,amsfonts,amssymb}
\usepackage{algorithm}
\usepackage{algorithmic}    % IEEEtran recommends algorithmic (not algorithmicx)
\usepackage{array}
\usepackage{bm}
\usepackage{booktabs}
\usepackage{multirow}
\usepackage{graphicx}
\usepackage{url}
\usepackage{cite}           % IEEE-style numeric citations
\usepackage{textcomp}
\usepackage{stfloats}       % allows figure* at bottom
% \usepackage[caption=false,font=normalsize,labelfont=sf,textfont=sf]{subfig}
\usepackage[caption=false,font=footnotesize,labelfont=sf,textfont=sf]{subfig}
\usepackage{siunitx}

% theorem/lemma environments
\newtheorem{theorem}{Theorem}
\newtheorem{lemma}{Lemma}

\hyphenation{op-tical net-works semi-conduc-tor IEEE-Xplore}

% ----------------------------------------------------------------
\begin{document}

% ---- Title ----
\title{Capacity-Constrained Target Assignment and Recurrent Multi-Agent Reinforcement Learning for Cooperative Interception of Maneuvering UAV Swarms}

% ---- Authors ----
% \author{Shuyi~Gao,~Defu~Lin,~Duo~Zheng,~and~Zhichen~Chu%
% \thanks{S. Gao, D. Lin, D. Zheng, and Z. Chu are with the School of Aerospace
% Engineering, Beijing Institute of Technology, Beijing 100081, China
% (e-mail: zhengduohello@126.com).}%
% \thanks{Manuscript received 2025.}}

%----双盲投稿--
\author{}
\markboth{Under review}%
{Anonymous Submission}

% ---- Running heads ----
\markboth{IEEE Transactions on Automation Science and Engineering}%
{Gao \MakeLowercase{\textit{et al.}}: MARL-Based Cooperative Interception for Maneuvering Swarm Targets}

\maketitle

% ---- Abstract ----
\begin{abstract}
The swarm attack of unmanned aerial vehicles (UAVs) poses a significant threat to future
battlefield operations.
Achieving efficient and precise cooperative interception against attacking UAV swarms with
highly dynamic maneuvering capabilities and adaptive penetration strategies poses a critical
challenge for UAV interceptor swarms.
This study addresses the cooperative interception of maneuvering targets in swarm attacks by
decomposing the problem into two key tasks: target assignment and cooperative interception
guidance.
For target assignment, an Improved Dung Beetle Optimization (IDBO) algorithm is proposed to
search directly over discrete assignments while enforcing target coverage and capacity.
For cooperative interception guidance, a recurrent Multi-Agent Proximal Policy Optimization
(R-MAPPO) controller is used under centralized training and decentralized execution. Its
three-dimensional observation and reward explicitly represent relative altitude, line-of-sight
geometry, group time-to-go dispersion, terminal interception, and control effort. A stored
capacity-feasible assignment artifact is passed unchanged to the guidance environment,
providing an auditable assignment-to-guidance interface rather than a regenerated end-to-end
optimizer run. Representative three-dimensional engagements, a 30-seed
assignment study, and separate software channel-impairment tests quantify both successful
operation and a delayed-coordination failure. The results support feasibility under the tested
conditions, while no claim of global optimizer convergence, Nash equilibrium, or hardware-in-
the-loop validation is made.
\end{abstract}

{\linespread{0.9}\selectfont
\hspace*{0.0em}\textbf{\textit{Note to Practitioners—}}\textbf{This paper is motivated by practical UAV swarm defense scenarios, where multiple interceptors must cooperatively engage maneuvering aerial targets under communication and onboard maneuver constraints. The revised framework exposes the interface between target assignment and guidance: a stored capacity-feasible assignment artifact is consumed unchanged by a recurrent MAPPO controller. The simulation now uses three-dimensional fixed-wing kinematics and reports representative trajectories together with separate tests of sensing delay, noise, command lag, and packet loss. These are mathematical and software channel simulations, not flight or hardware-in-the-loop experiments. The observed delayed-coordination case also shows that communication impairment can preserve individual interception while breaking group simultaneity; operational use therefore requires online health monitoring, reassignment, and hardware validation beyond the present study.}\par\vspace{0.5\baselineskip}
}

% ---- Keywords ----
\begin{IEEEkeywords}
Cooperative interception, deep reinforcement learning, centralized training with decentralized
execution, target assignment, UAV swarm.
\end{IEEEkeywords}


%% ===========================================================
\section{Introduction}
\label{sec:introduction}
\IEEEPARstart{A}{s} warfare evolves toward intelligent swarm operations, the likelihood of
encountering attacks from adversarial UAV swarms significantly increases during combat.
Consequently, the efficient interception of incoming swarm targets becomes a critical
imperative for future battlefield operations.
To address the challenge of cooperative interception against maneuvering attacking swarms with
adaptive penetration strategies, there is an urgent need to advance academic research in swarm
target interception technologies.
This study aims to enhance interceptors with intelligent cooperative interception capabilities,
enabling effective navigation of complex challenges in modern battlefield environments.

Multi-target cooperative interception encompasses two key tasks: target assignment and
cooperative interception.
The effectiveness of multi-target interception hinges on the optimality of target assignment,
establishing it as a primary focus in swarm target interception strategies.
As the number of UAVs increases, the complexity of target assignment escalates significantly,
resulting in a multi-constrained, multi-parameter optimization problem.
The objective is to develop a rapid and efficient target assignment method for allocating
specific attacking targets to interceptors in operational scenarios~\cite{KLINE2019226}.
The target assignment problem involves complex constraints, and the non-dominated sorting
algorithm can dynamically generate a set of reference points with good convergence and
distribution performance~\cite{WANG2025952,bai2024hybrid,zou2024moea}.
Moreover, employing an online operator selection mechanism can enhance the efficiency of the
algorithm~\cite{ahmad2024fuzzy,ming2024constrained}.
The target assignment problem can be solved using multi-information particle swarm optimization
to find the solution that yields the highest gain~\cite{mishra2015weapon,wei2020particle}.
Introducing information from surrounding particles can enhance the stability of the traditional
particle swarm algorithm while also addressing the issue of local convergence in the target
assignment problem~\cite{long2024improved}.
The genetic algorithm is a common approach for solving the target assignment
problem~\cite{he4914655modelling,li2020multi,sun2022multi,chang2018solving}, and a GA with a
novel crossover operator can further enhance algorithmic efficiency~\cite{weng2024diversity}.
Ant colony optimization is also an effective strategy for solving the target assignment
problem~\cite{li2024comprehensive}, and its combination with a greedy algorithm can improve
the accuracy of parallel computation~\cite{yi2024solving}.
Under multiple constraints, researchers have proposed various heuristic algorithms to solve the
target assignment problem, such as neural networks~\cite{li2024comprehensive}, auction
algorithm~\cite{wang2024collaborative}, tabu search~\cite{yao2024tabu}, simulated annealing
algorithm~\cite{khurshid2024hybrid}, VLSN~\cite{yue2024graph}.
However, in cooperative interception scenarios, the aforementioned algorithms struggle to
handle the highly dynamic nature of the battlefield due to their high time complexity, and they
fail to meet the dual requirements of both speed and optimality.

The effectiveness of the terminal guidance phase in multi-target interception is significantly
influenced by the cooperative guidance law, a critical factor requiring thorough consideration.
The target assignment strategy transforms swarm operations into a coordinated multi-to-one
interception problem, with current research focusing on three key areas: angular coordination,
temporal coordination, and intelligent cooperation.
Significant progress has been made in temporal coordination.
Classical proportional navigation-based guidance laws enable synchronized multi-missile attacks
on a single target~\cite{tao2024analytical}.
Adaptive control methods adjust attack timing based on real-time operational conditions, making
them suitable for intercepting maneuvering targets~\cite{wang2025generalized,zhang2022multiple}.
Sliding Mode Control (SMC) is widely adopted for cooperative
interception~\cite{sinha2020supertwisting}, with SMC incorporating miss distance and attack
time constraints enhancing interception
accuracy~\cite{shen2025three,jin2025cooperative,hu2019sliding}.
Considerable research has also been conducted on angular coordination.
From a traditional control perspective, backstepping control~\cite{li2024adaptive} and
leader-follower cooperative methods~\cite{he2017impact} significantly improve interception
efficiency.
SMC with angular constraints is prevalent in cooperative interception, encompassing approaches
based on impact angle constraints~\cite{wang2025hybrid} and line-of-sight angle
constraints~\cite{li2024near}.
Optimal control methods applied to impact angle-constrained cooperative guidance enhance
precision by integrating impact angle and line-of-sight angle
constraints~\cite{wang2025adaptive,tao2023optimal}.

Recent advancements in intelligent technologies have driven extensive research into cooperative
interception strategies for UAV swarms, with intelligent algorithms demonstrating substantial
potential in addressing complex multi-target scenarios~\cite{chen2024enhanced,wu2020collaborative}.
Optimization-based approaches, such as Particle Swarm Optimization (PSO), excel in
collaborative attack problems by autonomously optimizing resource allocation and maximizing
operational efficiency~\cite{zhao2024hierarchical,wu2020collaborative}.
Search-based algorithms, like A*~\cite{liao2024research}, enable efficient path planning for
coordinated attacks.
Control-based methods, including Model Predictive Control
(MPC)~\cite{kim2024adaptive}, provide robust guidance under dynamic constraints.
Moreover, the integration of neural networks with reinforcement learning, particularly deep
reinforcement learning (DRL), enhances the adaptability and intelligence of interception
strategies~\cite{luo2022uav,zhan2022multiple,wang2021uav}.
DRL enables UAV swarms to learn complex policies from high-dimensional data, optimizing target
assignment and maneuvering in real-time to counter maneuvering targets~\cite{sun2024memory}.
In contrast, conventional cooperative interception guidance methods rely heavily on
linearization and simplifying assumptions to design feasible guidance laws~\cite{zhao2025global}.
These methods often struggle to address the complexity of dynamic, non-linear environments
involving maneuvering swarm targets.
Data-driven intelligent cooperative interception methods, by contrast, leverage extensive
datasets to learn adaptive strategies, significantly reducing dependence on restrictive
assumptions~\cite{fainkich2025cooperative}.
By modeling complex interactions through neural networks and DRL, these methods achieve greater
robustness and flexibility.
Despite advancements, intelligent cooperative interception methods struggle with maneuvering
swarm targets in complex dynamic environments due to limited model generalization, low training
efficiency, and low interception success rates, particularly in large-scale swarm combat
requiring precise temporal coordination~\cite{horyna2023decentralized}.
Addressing these challenges requires further research into scalable, generalizable, and
efficient algorithms to fully realize the potential of intelligent cooperative interception.

In summary, existing cooperative interception methods for countering attacking UAV swarms with
adaptive penetration capabilities suffer from low interception success rates in complex dynamic
environments.
Specifically, target assignment methods exhibit low allocation efficiency, while cooperative
interception guidance strategies struggle with terminal overload saturation and insufficient
adaptability to dynamic target maneuvers in practical scenarios, hindering effective swarm
coordination.
Consequently, more research is necessary to address these challenges.

This study proposes a swarm cooperative interception method based on multi-agent reinforcement
learning and IDBO algorithm, enhancing target assignment efficiency and guidance strategy
performance to achieve precise and efficient interception of maneuvering targets.
The proposed intelligent cooperative interception algorithm offers the following advantages.

\begin{enumerate}
\item A capacity-constrained IDBO assignment method is formulated for the many-to-one
interception problem. Its discrete rolling, dancing, breeding, and stealing updates are
combined with feasibility repair and elitist acceptance. A 30-seed coefficient study and a
bounded-delay information-propagation experiment characterize empirical stability, scale, and
communication cost without claiming global or game-theoretic convergence.

\item A recurrent MAPPO guidance controller is specified consistently with the deployed
checkpoints. A 17-dimensional local observation captures three-dimensional relative geometry,
closing motion, group time-to-go dispersion, and previous commands; a GRU provides temporal
memory, and the actor outputs bounded axial, yaw-normal, and pitch-normal loads.

\item The assignment-to-guidance data interface is audited for 20 defenders and eight
targets using a stored capacity-feasible artifact; this is not presented as a regenerated
IDBO-to-guidance run. The revised evaluation reports representative three-dimensional R-MAPPO and
capacity-matched PN engagements, plus separate delay/noise/lag and software channel-
impairment tests. Reward shaping and PPO clipping are treated as standard task-adaptation and
training components rather than standalone methodological innovations.
\end{enumerate}

The paper is structured as follows: Section~\ref{sec:method} presents relevant problem modeling and foundational theory. Section~\ref{sec:target_assignment} and Section~\ref{sec:cooperative_algorithm} provide a detailed description of our approach, including its framework, modules, algorithms, and rationale. Section~\ref{sec:simulation} covers experiments, results, and discussions. Lastly, Section~\ref{sec:conclusion} summarizes the paper and outlines future directions.

%% ===========================================================
\section{Backgrounds and Preliminaries}
\label{sec:method}

In this section, we first introduce the modeling of multi-target interception of UAVs, and
then review reinforcement learning methods.

\subsection{Problem Modeling}

The multi-target intelligent interception problem can be decomposed into two subproblems:
target assignment and intelligent cooperative interception, with the swarm combat scenario
illustrated in Fig.~\ref{fenpei}.
In this scenario, the attacking swarm aims to penetrate a high-value target area, while the
defending swarm seeks to fully intercept the incoming UAVs.
During the target assignment phase, the defending swarm assesses its own flight status and
that of the attacking swarm using onboard sensors, subsequently allocating each defender to an
appropriate interception target via the target assignment strategy.
In the intelligent cooperative interception phase, defenders assigned to the same target employ
coordinated maneuvering strategies to intercept and neutralize the target.

\begin{figure}[!t]
\centering
\includegraphics[width=2.8in]{fenpei.pdf}
\caption{Cluster target interception problem.}
\label{fenpei}
\end{figure}

The object of study is a fixed-wing UAV. The horizontal engagement geometry in
Fig.~\ref{jianmo} is retained to define range and bearing, while the revised simulation and
guidance model propagate the full three-dimensional position, pitch, and yaw states.

\begin{figure}[!t]
\centering
\includegraphics[width=2.6in]{pic02.pdf}
\caption{Schematic diagram of confrontation and interception.}
\label{jianmo}
\end{figure}

The three-dimensional point-mass kinematic model is
\begin{equation}
\label{yundongxue}
\left\{\begin{aligned}
\dot{x}&=V\cos\vartheta\cos\psi,&
\dot{y}&=V\cos\vartheta\sin\psi,\\
\dot{z}&=V\sin\vartheta,&
\dot{V}&=g n_x,\\
\dot{\psi}&=\dfrac{g n_y}{V\max(\cos\vartheta,0.25)},&
\dot{\vartheta}&=\dfrac{g n_z}{V},
\end{aligned}\right.
\end{equation}
where $\mathbf{p}=[x,y,z]^\top$, $V$, $\psi$, and $\vartheta$ denote position, speed, yaw,
and pitch. The normalized commands $n_x$, $n_y$, and $n_z$ are axial, yaw-normal, and
pitch-normal loads, respectively.

In an interception mission, the relative motion relationship between the attacking and
defending UAVs can be described as:
\begin{equation}
\label{gongfang}
\left\{\begin{matrix}\dot{r} =V_{A}\cos(\gamma_{A}-q)-V_{D}\cos(\gamma_{D}-q) \\
r\dot{q} =V_{A}\sin(\gamma_{A}-q)-V_{D}\sin(\gamma_{D}-q)\end{matrix}\right.
\end{equation}
where $r$ is the inter-UAV distance, $q$ is the line-of-sight angle, $V_{A}$ ($V_{D}$) is the
velocity of the attacking (defending) UAV, and $\gamma_{A}$ ($\gamma_{D}$) is the
corresponding heading angle.

For the horizontal special case $\vartheta=n_z=0$, \eqref{yundongxue} reduces to the
two-dimensional relative-motion relations in \eqref{gongfang}. This reduction is used only
for defining the assignment geometry; all revised case simulations integrate
\eqref{yundongxue}.

During cooperative interception, the remaining impact time is defined as:
\begin{equation}
\label{shijian}
t_{go}=\frac{r}{\left|\dot{r}\right|}
\end{equation}

The Zero-Effort Miss (ZEM) is defined as:
\begin{equation}
\label{ZEM}
ZEM = \frac{r^{2}\cdot\dot{q}}{\left|\dot{r}\right|}
\end{equation}

\subsection{Reinforcement Learning}

Reinforcement learning (RL) is built on the Markov Decision Process (MDP), comprising state
space $\mathcal{S}$, action space $\mathcal{A}$, reward function $R$, policy $\pi$, and state
transition function $P$.
At each step, the agent selects action $a$ from state $s$, receives reward $r$, and
transitions to $s'$.
The key value functions are:
\begin{equation}
\label{RLV}
V^{\pi}(s)=E_{\pi}\left[\sum_{t=0}^{\infty}\gamma_{\mathrm{RL}}^{t}r_{t}\mid s_{0}=s\right]
\end{equation}
\begin{equation}
\label{RLQ}
Q^{\pi}(s,a)=E_{\pi}\left[\sum_{t=0}^{\infty}\gamma_{\mathrm{RL}}^{t}r_{t}\mid s_{0}=s,a_{0}=a\right]
\end{equation}
where $\gamma_{\mathrm{RL}}\in[0,1]$ is the discount factor.






\section{Target Assignment via Capacity-Constrained IDBO}
\label{sec:target_assignment}

This section formulates cooperative target assignment as a weighted bipartite
matching problem and solves it with a capacity-constrained Improved Dung Beetle
Optimization (IDBO) algorithm.

\subsection{Problem Formulation}
\label{subsec:problem_formulation}

Let $\mathcal{U} = \{u_1, u_2, \dots, u_M\}$ and $\mathcal{T} = \{t_1, t_2, \dots, t_N\}$ denote the interceptor and target sets, respectively. Figure \ref{km} illustrates the bipartite assignment structure.

\begin{figure}[!t]
\centering
\includegraphics[width=2.5in]{KM1.pdf}
\caption{Schematic diagram of weighted dichotomy.}
\label{km}
\end{figure}

The binary decision variable $X_{ij}=1$ if UAV $u_i$ is assigned to target $t_j$ and 0 otherwise. The edge weight $P_{ij}$ quantifies the interception probability. Intuitively, a desirable interceptor-target edge should have both a small predicted miss tendency and a favorable heading geometry; therefore, the weight combines ZEM-based feasibility and angular advantage:

\begin{equation}
\label{eq:interception_probability}
P_{ij} = w_{\mathrm{ZEM}} P_{\text{ZEM}}(i,j)
+ w_{\sigma} P_{\sigma}(i,j)
\end{equation}
where $w_{\mathrm{ZEM}}+w_{\sigma}=1$ and both weights are nonnegative.
The Zero-Effort Miss (ZEM) component is:

\begin{equation}
\label{eq:zem_prob}
P_{\text{ZEM}}(i,j) = \exp\left[-\frac{1}{2}\left(\frac{\text{ZEM}(i,j)}{\overline{\text{ZEM}}}\right)^2\right]
\end{equation}
where $\overline{\text{ZEM}} = \frac{1}{MN}\sum_{i=1}^{M}\sum_{j=1}^{N} \text{ZEM}(i,j)$.

The angular advantage component is:

\begin{equation}
\label{eq:angular_prob}
P_{\sigma}(i,j) = \exp\left[-\frac{1}{2}\left(\frac{S_{\sigma}(i,j)}{\overline{S_{\sigma}}}\right)^2\right]
\end{equation}
where $\overline{S_{\sigma}} = \frac{1}{MN}\sum_{i=1}^{M}\sum_{j=1}^{N} S_{\sigma}(i,j)$, with $S_{\sigma}(i,j) = |\gamma_{ij} - q_{ij}|$ denoting the heading error. The assignment objective below converts the tactical goal of ``neutralize as many attackers as possible'' into a survival-probability minimization problem. Under many-to-one assignment, each additional capable interceptor assigned to the same target multiplies that target's survival probability by another failure term:

\begin{equation}
\label{eq:optimization_formulation}
\begin{aligned}
\min_{\mathbf{X}} \quad & J = \sum_{j=1}^{N} \left[ \prod_{i=1}^{M} \left(1 - P_{ij}\right)^{X_{ij}} \right] \\
\text{s.t.} \quad & \sum_{j=1}^{N} X_{ij} = 1, \quad \forall i \in \{1,\dots,M\} \\
& \sum_{i=1}^{M} X_{ij} \geq 1, \quad \forall j \in \{1,\dots,N\} \\
& \sum_{i=1}^{M} X_{ij} \leq L_{\max}, \quad \forall j \in \{1,\dots,N\} \\
& X_{ij} \in \{0,1\}, \quad \forall i,j
\end{aligned}
\end{equation}
where $\prod_{i=1}^{M} (1 - P_{ij})^{X_{ij}}$ yields the survival probability of $t_j$, which decreases geometrically with the number of assigned interceptors.

\subsection{Capacity-Constrained IDBO and Input Dissemination}
\label{subsec:idbo_algorithm}

The optimizer uses a discrete assignment vector
\(\mathbf{x}=[x_1,\ldots,x_M]^\top\), where \(x_i=j\) assigns interceptor
\(u_i\) to target \(t_j\). This encoding implements the row constraint in
\eqref{eq:optimization_formulation} directly. A deterministic repair operator
\(\mathcal{R}\) moves the least-cost interceptor whenever a target is empty or
contains more than \(L_{\max}\) members; consequently, every evaluated candidate
satisfies all three capacity constraints.

At iteration \(k\), the population is divided equally among the four DBO roles.
The discrete updates have the following operational interpretation: rolling
copies selected components from the current best solution while retaining
occasional random moves; dancing randomly redirects selected components;
breeding recombines the best and second-best assignments; and stealing copies
selected components from the best assignment. The number of modified components
is controlled by
\begin{align}
\label{eq:idbo_schedule}
s_k&=0.32\left(1-\frac{k}{K-1}\right)+0.04,\\
m_k&=\max(1,\operatorname{round}(s_k M)).
\end{align}
Thus, the search moves from broader exploration toward local exploitation but
never becomes deterministic. Every candidate is repaired and accepted only if
its objective in \eqref{eq:optimization_formulation} does not exceed that of its
parent. The best-so-far value is also retained. These two explicit implementation
rules replace the previously ambiguous continuous preference equations.

Separately from the optimizer, we evaluate dissemination of fixed-snapshot input
sources over a connected sparse graph. Nodes repeatedly forward the source
identifiers they know, and delayed or dropped messages are retransmitted in
subsequent rounds. This experiment measures only how long immutable information
needs to reach all nodes. It does not execute distributed IDBO, assignment
consensus, or failure-triggered reassignment.

The complete reproducible procedure is summarized in
Algorithm~\ref{alg:idbo}. The assignment exported to guidance is a length-\(M\)
integer vector rather than a separately reconstructed topology.

\begin{algorithm}[!t]
\caption{Capacity-Constrained IDBO and Assignment Export}
\label{alg:idbo}
\footnotesize
\begin{algorithmic}[1]
\REQUIRE Probability matrix \(\mathbf{P}\), population size \(P\), iteration
budget \(K\), capacity \(L_{\max}\)
\ENSURE Feasible assignment vector \(\mathbf{x}^{*}\)
\STATE Initialize \(P\) random assignments and repair each with \(\mathcal{R}\)
\STATE Evaluate the target-survival objective in
Eq.~\eqref{eq:optimization_formulation}
\FOR{\(k=0,\ldots,K-1\)}
    \STATE Compute \(s_k\) and \(m_k\) from Eq.~\eqref{eq:idbo_schedule}
    \STATE Rank the population and retain the best and second-best assignments
    \FOR{each rolling, dancing, breeding, or stealing candidate}
        \STATE Modify \(m_k\) assignment components using the role-specific rule
        \STATE Apply capacity repair \(\mathcal{R}\)
        \STATE Accept the candidate only if its cost does not increase
    \ENDFOR
    \STATE Retain the best-so-far feasible assignment
\ENDFOR
\STATE Export \(\mathbf{x}^{*}\) and use it unchanged in the guidance environment
\end{algorithmic}
\end{algorithm}

\subsection{Convergence Scope, Complexity, and Delay Model}
\label{subsec:convergence_equilibrium}

Elitist acceptance gives the limited but rigorous property
\begin{equation}
\label{eq:incumbent_monotone}
J_{\mathrm{best}}^{(k+1)}\le J_{\mathrm{best}}^{(k)},\qquad
J_{\mathrm{best}}^{(k)}\ge 0 .
\end{equation}
Hence the recorded incumbent sequence converges to a finite limit. This does
\emph{not} prove that the population reaches consensus, that the limit is
globally optimal, or that the assignment is a Nash equilibrium. The problem is
nonconvex and discrete, and the stochastic search can terminate at a local
plateau. We therefore report paired-seed empirical cost and stabilization
statistics rather than an equilibrium claim.

For a candidate evaluation and repair implemented by direct scanning, the
worst-case work is \(O(MN)\). A population of size \(P\) and \(K\) iterations
therefore costs \(O(KPMN)\); memory is \(O(PM+MN)\). These are upper bounds for
the implementation, not hardware-independent runtime guarantees. The separate
input-source propagation experiment records synchronous message rounds. On a
connected graph with bounded link delay, repeated delivery eventually
disseminates immutable source information; this observation is not a convergence
result for IDBO. With changing data, permanent link failure, or a disconnected
graph, even this limited conclusion no longer applies.

\section{Recurrent MAPPO for Three-Dimensional Cooperative Interception}
\label{sec:cooperative_algorithm}

The revised guidance description is restricted to the architecture represented by
the evaluated checkpoints. It is a recurrent MAPPO controller under centralized
training with decentralized execution (CTDE). In particular, the deployed actor
contains an MLP encoder, a GRU, and a Gaussian action head. The earlier
trust-modulated exploration and attention-residual claims are not retained because
they were not independently supported by the evaluated checkpoint or an ablation.

\subsection{Actor--Critic Architecture and PPO Update}

For interceptor \(i\), the encoder and recurrent state update are
\begin{align}
\mathbf{e}_{i,t}&=f_{\mathrm{MLP}}(\mathbf{o}_{i,t};\theta_e),\\
\mathbf{h}_{i,t}&=\operatorname{GRU}(\mathbf{e}_{i,t},
\mathbf{h}_{i,t-1};\theta_h).
\label{eq:rmapppo_gru}
\end{align}
The actor defines a diagonal Gaussian and then clips the sampled command to the
flight envelope:
\begin{align}
\pi_\theta(\mathbf{a}_{i,t}\mid\mathbf{o}_{i,t},\mathbf{h}_{i,t-1})
 &=\mathcal{N}\!\left(\boldsymbol{\mu}_\theta(\mathbf{h}_{i,t}),
 \operatorname{diag}(\boldsymbol{\sigma}_\theta^2)\right),\nonumber\\[-2pt]
\mathbf{a}_{i,t}&=[n_{x,i},n_{y,i},n_{z,i}]^\top,\nonumber\\
\mathcal A&=[-0.1,1]\times[-1,1]\times[-1,1].
\label{eq:action3d}
\end{align}
The centralized critic receives the joint state during training, whereas each
actor uses only its local observation and recurrent state during execution:
\begin{equation}
s_t=[\mathbf{o}_{1,t}\|\cdots\|\mathbf{o}_{M,t}],\qquad
V_\phi(s_t)\in\mathbb{R}.
\end{equation}
Generalized advantage estimation is
\begin{align}
\delta_t&=r_t+\gamma_{\mathrm{RL}}V_\phi(s_{t+1})-V_\phi(s_t),\\
\hat A_t&=\sum_{\ell=0}^{T-t-1}
(\gamma_{\mathrm{RL}}\lambda_{\mathrm{GAE}})^\ell\delta_{t+\ell}.
\end{align}
The actor uses the standard clipped PPO surrogate
\begin{equation}
L^{\mathrm{clip}}(\theta)=
\mathbb{E}_t\!\left[
\min\!\left(r_t(\theta)\hat A_t,
\operatorname{clip}(r_t(\theta),1-\epsilon,1+\epsilon)\hat A_t\right)
\right],
\end{equation}
where \(r_t(\theta)=\pi_\theta(\mathbf a_t\mid\mathbf o_t)/
\pi_{\theta_{\mathrm{old}}}(\mathbf a_t\mid\mathbf o_t)\).
Clipping regularizes sampled policy updates but does not impose a hard bound on
the realized probability ratio and does not guarantee global convergence.

\begin{figure*}[!t]
\centering
\resizebox{0.97\textwidth}{!}{%
\mbox{%
\fbox{\parbox[c][0.62in][c]{1.18in}{\centering Local 3-D\\observation \(\mathbf{o}_{i,t}\)}}
\(\longrightarrow\)
\fbox{\parbox[c][0.62in][c]{1.12in}{\centering MLP feature\\\(\mathbf{e}_{i,t}\)}}
\(\longrightarrow\)
\fbox{\parbox[c][0.62in][c]{1.18in}{\centering GRU update\\
\(\mathbf{h}_{i,t-1}\!\to\!\mathbf{h}_{i,t}\)}}
\(\longrightarrow\)
\fbox{\parbox[c][0.62in][c]{1.15in}{\centering Gaussian\\actor head}}
\(\longrightarrow\)
\fbox{\parbox[c][0.62in][c]{1.20in}{\centering Bounded loads\\
\([n_x,n_y,n_z]\)}}}}\\[5pt]
\fbox{\parbox[c][0.54in][c]{2.05in}{\centering Joint state \(s_t\) during training}}
\(\longrightarrow\)
\fbox{\parbox[c][0.54in][c]{1.45in}{\centering Centralized critic \(V_\phi\)}}
\(\longrightarrow\)
\fbox{\parbox[c][0.54in][c]{1.40in}{\centering GAE and PPO update}}
\caption{Revised R-MAPPO data flow. The actor path shows one unambiguous temporal
transition from \(\mathbf{h}_{i,t-1}\) to \(\mathbf{h}_{i,t}\); the critic and
optimization path is used only during centralized training.}
\label{fig:algorithm_flow}
\end{figure*}

\subsection{Three-Dimensional Observation and Reward}
\label{subsec:state_reward}

The 17-dimensional local observation is
\begin{equation}
\label{eq:observation_space}
\begin{split}
\mathbf{o}_{i,t}=[&
r/2000,\ \Delta z/1000,\ \dot r/100,\ e_\angle/\pi,\
V_i/65,\ V_T/65,\\
&\vartheta_i/\pi,\ \psi_i/\pi,\ \boldsymbol{\ell}_{iT}^\top,\
(t_{go,i}-\bar t_{go})/50,\\
&\Delta t_{go}/50,\ (t_{go,i}-t_{go,\min})/50,\
n_{x,i}^{-},n_{y,i}^{-},n_{z,i}^{-}]^\top ,
\end{split}
\end{equation}
where \(\boldsymbol{\ell}_{iT}\in\mathbb{R}^3\) is the unit line-of-sight
vector, \(e_\angle\) is the angle between the interceptor velocity and the LOS,
\(\Delta t_{go}=t_{go,\max}-t_{go,\min}\) within the assigned group, and the
superscript \({}^{-}\) denotes the previous applied load. The observation uses
relative quantities and normalized scales so that the same actor can execute on
homogeneous interceptors.

The implemented reward has six task terms:
\begin{equation}
\label{eq:reward_function}
r_{i,t}=\sum_{q\in\{d,a,h,c,s,e\}}w_q r_{i,t}^{q},
\end{equation}
\begin{align}
r^d&=\exp(-\alpha_d r),&
r^a&=-\alpha_a e_\angle^{\,p_a},\\
r^h&=R_{\mathrm{hit}}\mathbb I(r<d_{\mathrm{hit}}),&
r^c&=-\alpha_c|t_{go,i}-\bar t_{go}|^{p_c}
+B_c,\\
r^s&=-\alpha_s(\Delta t_{go})^{p_s}+B_s,&
r^e&=-\alpha_e\|\mathbf a_{i,t}\|_2^2\Delta t .
\end{align}
The bonuses \(B_c\) and \(B_s\) are applied only inside their respective
coordination tolerances. The nominal implementation uses
\begin{align*}
(w_d,w_a,w_h,w_c,w_s,w_e)&=(0.1,1,1,1,1,1),\\
(\alpha_d,\alpha_a,\alpha_c)&=(10^{-3},10^{-2},5{\times}10^{-3}),\\
(\alpha_s,\alpha_e)&=(2{\times}10^{-2},5{\times}10^{-2}),\\
R_{\mathrm{hit}}&=3,\qquad \tau_{\mathrm{sync}}=0.5~\mathrm{s}.
\end{align*}

\subsection{Sensitivity and Convergence Scope}
\label{subsec:reward_sensitivity}

Reward sensitivity is best interpreted through invariant mission metrics rather
than returns, because changing a reward weight changes the numerical definition
of return. Increasing \(w_d\) favors early range reduction; increasing \(w_a\)
favors velocity--LOS alignment; increasing \(w_c\) or \(w_s\) prioritizes
within-group timing; and increasing \(w_e\) suppresses load usage. Excessive
coordination weight can delay a geometrically favorable interceptor, whereas
excessive energy weight can prevent the maneuver needed for a hit. Accordingly,
the evaluation reports hit count, group arrival spread, miss distance, and load
limits rather than comparing returns across reward definitions. A full
multi-seed retraining sweep of these weights is not claimed by the present
evidence; this limitation is stated explicitly to avoid interpreting an
uncontrolled one-seed preview as sensitivity evidence.

R-MAPPO and conventional MAPPO share the same PPO-style nonconvex optimization
and therefore the same lack of a global or Nash-convergence guarantee. The GRU
changes the policy class by retaining observation history, which is useful under
partial observation, but it does not strengthen the theoretical convergence
guarantee. PPO clipping is a training-time operation and adds no actor operations
during deployment. The deployed actor contains 533,800 parameters; inference
uses only the MLP, GRU, and Gaussian mean head. Processor-specific latency was
not measured, so no onboard deadline claim is made.

\subsection{Evaluation Criteria}
\label{subsec:eval}

For each target group \(\mathcal G_j\), the terminal synchronization spread and
miss distance are
\begin{align}
E_{\mathrm{sync},j}&=\max_{i\in\mathcal G_j}t_{f,i}
-\min_{i\in\mathcal G_j}t_{f,i},\\
E_{\mathrm{miss}}&=\frac{1}{M}\sum_{i=1}^{M}
\|\mathbf p_i(t_{f,i})-\mathbf p_{T(i)}(t_{f,i})\|_2 .
\end{align}
We additionally report the number of interceptor hits, the number of covered and
synchronized target groups, and the maximum absolute command in each of the
three load channels. A group is synchronized in the final case study when
\(E_{\mathrm{sync},j}\le0.5\) s.

\begin{algorithm}[!t]
\caption{R-MAPPO Cooperative Interception Guidance}
\label{alg:r_mappo}
\footnotesize
\begin{algorithmic}[1]
\REQUIRE Exported assignment \(\mathbf{x}^{*}\), 3-D environment, horizon \(T\)
\ENSURE Decentralized recurrent actor \(\pi_\theta\)
\STATE Form target groups directly from \(\mathbf{x}^{*}\)
\FOR{each training episode}
    \STATE Reset 3-D states and recurrent states
    \FOR{\(t=0,\ldots,T-1\)}
        \FOR{each interceptor \(i\)}
            \STATE Build \(\mathbf{o}_{i,t}\) from Eq.~\eqref{eq:observation_space}
            \STATE Update \(\mathbf{h}_{i,t}\) with Eq.~\eqref{eq:rmapppo_gru}
            \STATE Sample and clip \([n_x,n_y,n_z]\) using Eq.~\eqref{eq:action3d}
        \ENDFOR
        \STATE Propagate 3-D dynamics and compute Eq.~\eqref{eq:reward_function}
        \STATE Store local observations, joint state, actions, rewards, and masks
    \ENDFOR
    \STATE Compute GAE and update actor/critic with clipped PPO
\ENDFOR
\STATE During deployment, execute each actor from local observation and GRU state
\end{algorithmic}
\end{algorithm}


\section{Simulation Verification}
\label{sec:simulation}

\subsection{Revised Three-Dimensional Setup}
\label{subsec:setup}

The revised evaluation uses 20 defending UAVs and eight attacking UAVs with the
three-dimensional dynamics in \eqref{yundongxue}. The simulation time step is
0.05 s and the maximum horizon is 1500 steps (75 s), correcting the previously
reported 25-step value, which belonged to an unrelated auxiliary benchmark.
Table~\ref{tab:revised_setup} gives the settings used by the saved case presets.

\begin{table}[!t]
\centering
\caption{Revised Three-Dimensional Evaluation Settings}
\label{tab:revised_setup}
\footnotesize
\renewcommand{\arraystretch}{1.12}
\begin{tabular}{p{0.52\columnwidth}p{0.36\columnwidth}}
\toprule
Setting & Value \\
\midrule
Defenders / attackers & 20 / 8 \\
Time step / maximum horizon & 0.05 s / 1500 steps \\
Initial defender radius & 52.63--65.44 m \\
Initial attacker radius & 1987.72--2145.95 m \\
Initial defender / attacker altitude & 0 m / 120 m \\
Initial defender speed & 20.32--30.02 m/s \\
Initial attacker speed & 30.23--45.87 m/s \\
Hit radius / synchronization threshold & 3 m / 0.5 s \\
Defender load limits & \(n_x\in[-0.1,1]\), \(n_y,n_z\in[-1,1]\) \\
\bottomrule
\end{tabular}
\end{table}

These explicit radial intervals replace the ambiguous \(D_1,D_2,D_3\) notation
and the inconsistent two-dimensional initialization sketch in the original
manuscript. The two attacker profiles are: Case 1, a piecewise evasive maneuver;
and Case 2, continuous sinusoidal weaving. Both include vertical motion. The
saved trajectories start at the audited states in Table~\ref{tab:revised_setup}
and use identical initial horizontal states for R-MAPPO and PN.

\subsection{IDBO Ablation, Weight Sensitivity, and Scale}
\label{subsec:idbo_results}

The revised assignment study implements the \(20\times8\) many-to-one decision
in \eqref{eq:optimization_formulation}, rather than the inverse eight-variable
surrogate used by the previous figures. Each variant uses 30 paired seeds, a
population of 40, and 100 iterations. All 60 final assignments in the coefficient
comparison satisfy the one-to-three interceptor capacity constraint.

\begin{figure}[!t]
\centering
\includegraphics[width=3.45in]{idbo_ablation_delay.pdf}
\caption{Capacity-constrained IDBO study. Left: best target-survival cost for a
common decaying mutation strength and a fixed strength (mean and one standard
deviation over 30 seeds). Right: rounds required for fixed-source information
propagation with 1\% packet dropout; this is not an IDBO convergence curve.}
\label{fig:idbo_revision}
\end{figure}

\begin{table}[!t]
\centering
\caption{IDBO Coefficient-Schedule Study (30 Paired Seeds)}
\label{tab:idbo_ablation}
\scriptsize
\setlength{\tabcolsep}{2.5pt}
\begin{tabular}{p{0.29\columnwidth}p{0.32\columnwidth}p{0.24\columnwidth}}
\toprule
Schedule & Final cost & Median iteration to 1\% band \\
\midrule
Decaying \(0.36\to0.04\) & \(0.006975\pm0.000141\) & 19.0 \\
Fixed \(0.18\) & \(0.006942\pm0.000083\) & 21.5 \\
\bottomrule
\end{tabular}
\end{table}

The decaying schedule reaches its own 1\%-of-final band slightly earlier, but it
does not improve final cost: among the 30 paired runs it is better, equal, and
worse in 1, 25, and 4 runs, respectively. We therefore do not claim a final-
quality advantage from coefficient adaptation. Likewise, varying the ZEM weight
\(w_{\mathrm{ZEM}}\in\{0.3,0.5,0.7\}\), with
\(w_\sigma=1-w_{\mathrm{ZEM}}\), changes the optimized objective mean to
0.004611, 0.006986, and 0.009956. These absolute values are not ranked because
the objective itself changes with the weight; the sweep demonstrates the
expected tactical tradeoff rather than a universal best weight.

For 60 iterations and population 30, the final archived audit recorded mean
runtimes of 0.417, 0.732, and 1.846 s for
\(M/N=20/8,40/16,80/32\), respectively. Repeated audit executions showed that
the absolute wall-clock values shift substantially with shared-host load;
accordingly, only the increasing scale trend and the analytical
\(O(KPMN)\) implementation bound are retained. These timings are not evidence
of an onboard deadline.
In the separate 20-node input-propagation test, the mean rounds at 1\% dropout
were 4.00, 5.80, 7.07, and 9.17 for maximum link delays of 0, 1, 2, and 4
rounds. Removing one node still allowed the 19-node connected subgraph to
disseminate all remaining source identifiers in \(7.17\pm0.46\) rounds, but no
task reassignment or guidance after vehicle loss was evaluated.

\subsection{Assignment-to-Guidance Interface}
\label{subsec:end_to_end}

Figure~\ref{fig:assignment_topology} shows the stored assignment artifact used
by the two guidance cases. Its one-based target vector is
\[
[1,2,3,4,5,6,7,8,\;1,2,3,4,5,6,7,8,\;1,2,3,4],
\]
so targets 1--4 receive three defenders and targets 5--8 receive two. A direct
reset audit confirms that the same vector is stored in the three-dimensional
case preset and consumed without remapping by the guidance environment. This
verifies the data interface between the stored assignment and guidance stages.
It does not claim that the representative assignment produced by the separate
30-seed reviewer study is identical to this legacy case artifact.

\begin{figure}[!t]
\centering
\includegraphics[width=3.2in]{target_assignment_topology.png}
\caption{Stored 20-to-8 capacity-feasible assignment passed unchanged to both
three-dimensional guidance cases.}
\label{fig:assignment_topology}
\end{figure}

\subsection{Representative Three-Dimensional Engagements}
\label{subsec:cases}

R-MAPPO and an \(N=4\) PN reference are evaluated from the audited initial
states. PN uses group-average time-to-go for longitudinal coordination and the
same axial and normal load limits. Figure~\ref{fig:case_3d_trajectories} presents
the saved representative trajectories.

\begin{figure*}[!t]
\centering
\subfloat[R-MAPPO, Case 1]{\includegraphics[width=3.25in]{rev_mappo_case1_3d.png}}
\hfill
\subfloat[PN, Case 1]{\includegraphics[width=3.25in]{rev_pn_case1_3d.png}}\\
\subfloat[R-MAPPO, Case 2]{\includegraphics[width=3.25in]{rev_mappo_case2_3d.png}}
\hfill
\subfloat[PN, Case 2]{\includegraphics[width=3.25in]{rev_pn_case2_3d.png}}
\caption{Representative three-dimensional engagements for the piecewise evasive
(Case 1) and sinusoidal weaving (Case 2) targets. Solid and dashed segments show
defender and target motion, respectively.}
\label{fig:case_3d_trajectories}
\end{figure*}

\begin{table*}[!t]
\centering
\caption{Representative-Run Three-Dimensional Results (\(n=1\) per Method/Case)}
\label{tab:case_results}
\footnotesize
\begin{tabular}{lcccccc}
\toprule
Method & Case & Interceptor hits & Groups hit & Groups synchronized &
Mean group spread (s) & Maximum spread (s) \\
\midrule
R-MAPPO & 1 & 20/20 & 8/8 & 8/8 & 0.1250 & 0.2500 \\
PN (\(N=4\)) & 1 & 20/20 & 8/8 & 8/8 & 0.1563 & 0.4000 \\
R-MAPPO & 2 & 20/20 & 8/8 & 8/8 & 0.1688 & 0.3000 \\
PN (\(N=4\)) & 2 & 20/20 & 8/8 & 8/8 & 0.1438 & 0.3500 \\
\bottomrule
\end{tabular}
\end{table*}

All saved impacts satisfy the 3-m criterion. The maximum R-MAPPO command
magnitudes \([|n_x|,|n_y|,|n_z|]\) are
\([0.9882,0.9975,0.6061]\) in Case 1 and
\([1.0000,1.0000,0.7749]\) in Case 2. Both methods succeed in these selected
episodes; therefore, Table~\ref{tab:case_results} supports feasibility and
describes synchronization, not statistical superiority or a success-rate
difference. Method-to-method interception times are not compared because the
two archived evaluators aggregate group time differently.

\begin{figure*}[!t]
\centering
\subfloat[R-MAPPO, Case 1]{\includegraphics[width=3.25in]{rev_mappo_case1_sync.png}}
\hfill
\subfloat[PN, Case 1]{\includegraphics[width=3.25in]{rev_pn_case1_sync.png}}\\
\subfloat[R-MAPPO, Case 2]{\includegraphics[width=3.25in]{rev_mappo_case2_sync.png}}
\hfill
\subfloat[PN, Case 2]{\includegraphics[width=3.25in]{rev_pn_case2_sync.png}}
\caption{Within-group terminal arrival-time spreads for the four representative
runs. The dashed threshold is 0.5 s.}
\label{fig:case_sync}
\end{figure*}

\subsection{Delay, Noise, Packet Loss, and Failure Case}
\label{subsec:robustness}

The final trajectories above contain no injected communication impairment. We
therefore report two separate tests and keep their criteria distinct. First, one
illustrative mathematical-simulation run per case applies a 50-ms sensing delay,
3-m position-noise standard deviation, 0.3-m/s velocity-noise standard deviation,
and a 0.30-s first-order command lag. Both cases record 20/20 hits and 8/8
synchronized groups with a 0.25-s maximum spread. This is a selected run, not a
robust success-rate estimate.

Second, a software channel-emulation test applies a 100-ms observation delay,
50-ms action delay, normalized observation/action noise standard deviations of
0.003/0.02, and 1\% command dropout to five policy endpoints. It uses randomized
geometry, a 12-m hit radius, and a 3-s synchronization threshold, so its rates
are not directly comparable with Table~\ref{tab:case_results}.

\begin{table}[!t]
\centering
\caption{Software Temporal-State and Channel Diagnostics (10 Episodes)}
\label{tab:channel_results}
\footnotesize
\begin{tabular}{llccc}
\toprule
Case & Diagnostic & All-hit & All-sync & Worst spread (s) \\
\midrule
1 & Full history & 10/10 & 10/10 & 0.60 \\
1 & GRU reset & 10/10 & 10/10 & 0.60 \\
1 & Channel impaired & 10/10 & 9/10 & 43.65 \\
2 & Full history & 10/10 & 10/10 & 0.20 \\
2 & GRU reset & 10/10 & 9/10 & 39.55 \\
2 & Channel impaired & 10/10 & 10/10 & 0.20 \\
\bottomrule
\end{tabular}
\end{table}

The full-history and GRU-reset rows use identical seeds and the same recurrent
checkpoint; the reset diagnostic zeros its hidden state before every action. Resetting
preserves all-hit outcomes, but Case-2 episode 7 synchronizes only seven of eight
groups and reaches a 39.55-s worst spread. This inference-only intervention supports
the relevance of temporal state in that executed case. It is not a retrained
feed-forward baseline or a multi-seed architectural ablation, so no independent
causal performance gain is claimed for the GRU.

The Case-1 impaired episode 6 is a delayed-coordination failure: all 20
interceptors reach their targets, but only seven of eight groups meet the 3-s
criterion and the worst group spread is 43.65 s. This distinction matters
operationally because individual hit success does not imply coordinated success.
The endpoint objects ran in one Python process; the experiment is software
channel emulation, not hardware-in-the-loop or semi-physical validation.

Partial interceptor loss, stuck actuators, and held-out maneuver families were
not executed under the final 3-m cases. The present evidence therefore does not
establish failure-reassignment robustness or out-of-distribution generalization.
Real deployment additionally requires timestamped network transport, estimator
bias/drift models, processor latency and jitter measurements, actuator dynamics,
collision avoidance, terrain/weather effects, and flight or genuine HIL tests.


\section{Conclusion}
\label{sec:conclusion}

This paper formulated capacity-constrained many-to-one target assignment and
three-dimensional cooperative guidance as two auditable stages. The revised IDBO
study uses the same \(20\times8\) survival-probability objective as the manuscript
and shows monotone incumbent cost under explicit elitist acceptance, while also
showing that the tested coefficient decay does not improve final cost over the
fixed setting. The recurrent MAPPO description now matches the evaluated actor
architecture and three-dimensional action space.

In the four representative case runs, both R-MAPPO and capacity-matched PN
intercept all targets, with maximum within-group spreads of 0.25/0.30 s and
0.40/0.35 s, respectively. A separate software channel test exposes one delayed
coordination failure despite complete individual interception. These results
support feasibility under the stated conditions, not global optimizer
convergence, policy equilibrium, statistical superiority, or hardware readiness.
Future work will require genuine distributed optimization, multi-seed reward and
architectural retraining, loss-triggered reassignment, held-out maneuver tests, estimator and
actuator models, and processor-in-the-loop, HIL, and flight validation.

% ---- Bibliography ----
\bibliographystyle{IEEEtran}
\bibliography{cja-template}

\appendices
\section{Notation Table}
\label{app:notation}

Table~\ref{tab:notation} summarizes the main symbols used in the target assignment and
R-MAPPO formulations.

\begin{table}[!t]
\centering
\caption{Concise Notation Summary}
\label{tab:notation}
\renewcommand{\arraystretch}{1.15}
\footnotesize
\begin{tabular}{p{0.28\linewidth}p{0.62\linewidth}}
\toprule
Symbol & Description \\
\midrule
$\mathcal{U},\mathcal{T}$ & Interceptor and target sets. \\
$M,N$ & Numbers of interceptors and targets. \\
$X_{ij}$, $\mathbf{X}$ & Binary assignment variable and assignment matrix. \\
$P_{ij}$ & Interception probability weight for pair $(u_i,t_j)$. \\
$J$ & Target-assignment objective function. \\
$L_{\max}$ & Maximum number of interceptors assigned to one target. \\
$\mathbf{x}$, $x_i$ & Discrete assignment vector and target selected by interceptor \(u_i\). \\
$K,P$ & IDBO iteration budget and population size. \\
$s_k,m_k$ & Mutation strength and number of modified components at iteration \(k\). \\
$w_{\mathrm{ZEM}},w_\sigma$ & ZEM and angular weights in \(P_{ij}\). \\
$\mathbf{o}_i^{(t)}$ & Local observation of interceptor $i$ at time step $t$. \\
$s_t$ & Centralized state used by the critic. \\
$\mathbf{e}_{i,t},\mathbf{h}_{i,t}$ & Encoded feature and GRU hidden state. \\
$\mathbf{a}_i=[n_x,n_y,n_z]^\top$ & Three-dimensional normalized load command. \\
$\pi_{\theta_i}$, $V_{\phi}$ & Decentralized actor policy and centralized value function. \\
$r_t(\theta)$, $\hat{A}_t$ & PPO probability ratio and advantage estimate. \\
$\gamma_{\mathrm{RL}},\lambda_{\mathrm{GAE}},\epsilon$ & Discount, GAE, and PPO-clipping parameters. \\
$r_i^{(t)}$ & Reward of interceptor $i$ at time step $t$. \\
$w_q,\alpha_q$ & Reward weights and term-specific shaping coefficients. \\
$d_{\mathrm{hit}},\tau_{\mathrm{sync}}$ & Hit radius and group synchronization threshold. \\
$\mathcal{G}_i$ & Interceptors assigned to the same target as $u_i$. \\
$E_{\mathrm{sync},j}$, $E_{\mathrm{miss}}$ & Group arrival spread and mean miss-distance metrics. \\
\bottomrule
\end{tabular}
\end{table}

\end{document}
